\chapter{Monotone Folgen und bestimmte Divergenz}
\label{b21guide:ch:monotone}

Die Zielmenge ist weiterhin $\mathbb R$. Die Zeichen
$+\infty$ und $-\infty$ beschreiben ein Verhalten der Folge;
sie werden nicht als zusätzliche reelle Zahlen eingeführt.

\section{Monotone Konvergenz}

\begin{definition}[Monotone reelle Folge]
\label{b21guide:def:monotone}
Eine Folge heißt \emph{monoton wachsend}, wenn $a_n\leq a_{n+1}$
für alle $n$ gilt, und \emph{monoton fallend}, wenn
$a_{n+1}\leq a_n$ für alle $n$ gilt.
Bei strikten Ungleichungen heißt sie \emph{strikt monoton}.
Sie heißt \emph{schließlich monoton}, wenn eine Restfolge monoton ist.
\end{definition}

\begin{theorem}[Monotone Konvergenz]
\label{b21guide:thm:monotone-convergence}
Eine monoton wachsende reelle Folge ist genau dann konvergent,
wenn sie nach oben beschränkt ist.
Eine monoton fallende reelle Folge ist genau dann konvergent,
wenn sie nach unten beschränkt ist.

Eine schließlich monotone reelle Folge ist somit genau dann
konvergent, wenn sie beschränkt ist.
\end{theorem}

\begin{theorem}[Monotone Einschließung mit verschwindender Lücke]
\label{b21guide:thm:nested}
Sei $(u_n)$ monoton wachsend, $(v_n)$ monoton fallend und
$u_n\leq v_n$ für alle $n$.
Dann existieren reelle Grenzwerte $U,V$ mit $U\leq V$.
Gilt zusätzlich $v_n-u_n\to0$, so ist $U=V$;
jede Folge mit $u_n\leq a_n\leq v_n$ schließlich
konvergiert dann gegen diesen gemeinsamen Wert.
\end{theorem}

\section{Divergenz gegen Unendlich}

\begin{definition}[Bestimmte Divergenz]
\label{b21guide:def:infinity}
Für eine reelle Folge bedeuten
\[
 \begin{aligned}
 a_n\to+\infty
 &\ \Longleftrightarrow\
 \forall C\in\mathbb R\;\exists N\in\mathbb N\;
 \forall n\geq N:\ a_n>C,\\
 a_n\to-\infty
 &\ \Longleftrightarrow\
 \forall C\in\mathbb R\;\exists N\in\mathbb N\;
 \forall n\geq N:\ a_n<C.
 \end{aligned}
\]
In diesen Fällen heißt die Folge \emph{bestimmt divergent}.
Eine Folge, die weder einen reellen Grenzwert besitzt noch
gegen $+\infty$ oder $-\infty$ strebt, heißt \emph{unbestimmt divergent}.
\end{definition}

\begin{theorem}[Vollständige Alternative für monotone Folgen]
\label{b21guide:thm:monotone-dichotomy}
Eine monoton wachsende reelle Folge konvergiert entweder gegen
eine reelle Zahl oder strebt gegen $+\infty$; der zweite Fall
tritt genau bei fehlender oberer Beschränktheit ein.
Entsprechend konvergiert eine monoton fallende Folge reell
oder strebt bei fehlender unterer Beschränktheit gegen $-\infty$.
Dasselbe gilt jeweils bei schließlicher Monotonie.
\end{theorem}

\begin{theorem}[Teilfolgen aus Unbeschränktheit]
\label{b21guide:thm:unbounded-subsequence}
Eine reelle Folge ist genau dann nach oben unbeschränkt,
wenn sie eine Teilfolge besitzt, die gegen $+\infty$ strebt.
Nach unten unbeschränkt ist sie genau dann, wenn eine Teilfolge
gegen $-\infty$ strebt.
Unbeschränktheit der ganzen Folge allein erzwingt keine bestimmte
Divergenz der ganzen Folge.
\end{theorem}

\begin{theorem}[Restfolgen, Teilfolgen und Vergleiche]
Ist $a_n\to+\infty$, so gilt dies auch für jede Teilfolge;
endlich viele Änderungen beeinflussen die bestimmte Divergenz nicht.
Aus $a_n\leq b_n$ schließlich und $a_n\to+\infty$ folgt
$b_n\to+\infty$.
Entsprechend folgt aus $b_n\leq a_n$ schließlich und
$a_n\to-\infty$, daß $b_n\to-\infty$.
\end{theorem}

\begin{theorem}[Rechenregeln für bestimmte Divergenz]
\label{b21guide:thm:infinite-rules}
Es gelte $a_n\to+\infty$. Dann gelten:
\begin{enumerate}
 \item $-a_n\to-\infty$; für jedes $c>0$ gilt $ca_n\to+\infty$.
 \item Ist $(b_n)$ nach unten beschränkt, so ist
 $a_n+b_n\to+\infty$.
 \item Gilt schließlich $b_n\geq c$ für ein festes $c>0$,
 so ist $a_nb_n\to+\infty$.
 \item Es gilt $1/a_n\to0$ auf einer Restfolge, auf der $a_n>0$ ist.
\end{enumerate}
Die entsprechenden Aussagen für negative Divergenz ergeben sich
durch Vorzeichenwechsel. Ausdrücke wie $\infty-\infty$ oder
$0\cdot\infty$ sind dadurch nicht definiert.
\end{theorem}

\begin{theorem}[Kehrwert einer einseitigen Nullfolge]
Gelte $a_n>0$ schließlich. Dann gilt auf dieser Restfolge
\[
 a_n\to0\quad\Longleftrightarrow\quad\frac1{a_n}\to+\infty.
\]
Für schließlich negative $a_n$ gilt entsprechend
$a_n\to0$ genau dann, wenn $1/a_n\to-\infty$.
\end{theorem}

\section{Divergenz ohne reelle Häufungswerte}

\begin{theorem}[Vollständiges Kriterium für fehlende Häufungswerte]
\label{b21guide:thm:no-cluster}
Für jede reelle Folge $(a_n)$ gilt
\[
 H(a)=\varnothing\quad\Longleftrightarrow\quad |a_n|\to+\infty.
\]
Die beiden Divergenzrichtungen sind dadurch charakterisiert:
\begin{enumerate}
 \item $a_n\to+\infty$ genau dann, wenn $H(a)=\varnothing$
 und $(a_n)$ nach unten beschränkt ist.
 \item $a_n\to-\infty$ genau dann, wenn $H(a)=\varnothing$
 und $(a_n)$ nach oben beschränkt ist.
\end{enumerate}
Das Fehlen reeller Häufungswerte allein legt somit noch keine
der beiden Divergenzrichtungen fest.
\end{theorem}
